\documentclass[conference]{IEEEtran}

\usepackage{amsmath}
\usepackage{booktabs}
\usepackage{float}
\usepackage{graphicx}
\usepackage{microtype}
\usepackage{siunitx}
\usepackage[hidelinks]{hyperref}

\hypersetup{
  pdftitle={Evaluation-Protocol Sensitivity of COVID-19 Respiratory-Audio Screening Models},
  pdfauthor={Nishant Harkut and Ritu Tiwari},
  pdfkeywords={COVID-19, respiratory audio, evaluation design, participant-disjoint validation, multimodal fusion}
}

\graphicspath{{figures/}}
\sisetup{detect-all,round-mode=places,round-precision=3}

\title{Evaluation-Protocol Sensitivity of COVID-19 Respiratory-Audio Screening Models}

\author{
\IEEEauthorblockN{Nishant Harkut and Ritu Tiwari}
\IEEEauthorblockA{Department of Information Technology\\
Atal Bihari Vajpayee Indian Institute of Information Technology and Management\\
Gwalior, India}
}

\begin{document}
\maketitle

\begin{abstract}
Audio-based COVID-19 screening models have reported high discrimination, but performance estimates vary with evaluation design. Han \emph{et al.} showed biased participant allocation can inflate AUROC from 0.71 to 0.90. Coppock \emph{et al.} found matching reduced audio AUROC from 0.85 to 0.62 and symptoms alone outperformed audio. These findings raise a question: what would the performance of a strong audio model be under rigorously isolated participant-disjoint evaluation? We analyzed 2,088 quality-passing Coswara participants with fixed train/validation/test partitions (N=1,460/312/316). Training-only feature ranking retained 800 candidates from 10,140 numeric descriptors. Validation-selected multimodal fusion achieved test AUROC 0.895 (95\% bootstrap CI 0.852--0.933). Speech alone achieved AUROC 0.888. The fusion estimate was 0.007 AUROC above speech, but the paired difference was not statistically significant. Metadata alone achieved AUROC 0.964; shuffled-label retraining produced AUROC 0.503. These results show strong internal discrimination is achievable under participant-disjoint evaluation, but does not establish external validity, temporal stability, or superiority over metadata baselines. The estimate applies only to the Coswara test cohort and does not generalize to other datasets or time periods.
\end{abstract}

\begin{IEEEkeywords}
COVID-19, respiratory audio, evaluation design, participant-disjoint validation, bias in machine learning, multimodal fusion
\end{IEEEkeywords}

\section{Introduction}

Audio-based screening for COVID-19 has been proposed as a low-cost triage tool, with reported AUROC values reaching 0.90--0.97 in some studies \cite{laguarta2020cough,bhattacharya2023coswara,aytekin2024hst}. These estimates, however, depend critically on how participants are partitioned and how model selection is performed.

Han \emph{et al.} demonstrated that biased participant allocation can change AUROC from 0.71 under realistic evaluation to 0.90 under biased splits, using the same model architecture \cite{han2022realistic}. Coppock \emph{et al.} found that matching on confounders reduced audio classifier AUROC from 0.85 to 0.62, and that simple symptom checkers outperformed audio-based models \cite{coppock2024audio}. Avila \emph{et al.} showed that non-nested feature selection can produce optimistic AUROC estimates compared to nested evaluation \cite{avila2021feature}. These studies indicate that evaluation protocol is a first-order determinant of reported performance.

This matters for deployment. If internal estimates are optimistic due to participant leakage, test-guided selection, or feature-selection bias, then deployment expectations may be unrealistic. However, the magnitude of this bias under participant-disjoint, validation-only selection has not been quantified for a comprehensive Coswara workflow.

We ask: under rigorously isolated participant-disjoint evaluation with validation-only model selection, what discrimination does a strong multimodal audio model achieve, and how does this compare to metadata baselines and prior Coswara results? The contributions are: (1) a complete participant-audited pipeline with training-only feature ranking and validation-only model selection; (2) test results with uncertainty quantification for both fusion and single-modality branches; (3) metadata-only and shuffled-label baselines that establish non-audio predictiveness and the absence of gross label leakage; and (4) comparison to prior internal Coswara estimates with explicit acknowledgment of protocol differences.

The study is deliberately narrow: it estimates internal performance under one fixed participant-disjoint protocol. It does not evaluate temporal stability, cross-dataset transfer, or clinical deployment readiness. These require separate evaluation protocols.

\section{Materials and Methods}

\subsection{Study Design and Participant Flow}

The indexed Coswara release contained 2,746 participants. We excluded 632 participants with unresolved labels and 26 participants failing recording-quality checks, yielding a modeling cohort of 2,088 participants with resolved positive or negative COVID-19 status (Figure~\ref{fig:flow}).

\begin{figure}[!t]
\centering
% FLOW DIAGRAM PLACEHOLDER - should show:
% 2,746 indexed → 2,114 resolved → 2,088 quality-passing
% → Train 1,460 / Val 312 / Test 316
\includegraphics[width=\linewidth]{figures/participant_flow.pdf}
\caption{Participant flow. Of 2,746 indexed participants, 2,088 quality-passing participants with resolved labels entered fixed train/validation/test partitions. Exclusions were 632 unresolved labels and 26 quality failures.}
\label{fig:flow}
\end{figure}

Participants were assigned to train, validation, or test partitions (70/15/15) using stratified sampling on label, age band, and gender where cell sizes permitted, with fallback to label-only stratification. All recordings from one participant remained in one partition. The partition seed was 42.

\begin{table}[!b]
\caption{Quality-Passing Participant-Disjoint Coswara Partitions}
\label{tab:cohort}
\centering
\footnotesize
\begin{tabular}{lrrrr}
\toprule
Partition & Participants & Positive & Negative & Prevalence \\
\midrule
Train      & 1,460 & 476 & 984 & 32.6\% \\
Validation &   312 & 101 & 211 & 32.4\% \\
Test       &   316 & 103 & 213 & 32.6\% \\
\midrule
\textbf{Total} & \textbf{2,088} & \textbf{680} & \textbf{1,408} & \textbf{32.6\%} \\
\bottomrule
\end{tabular}
\end{table}

\subsection{Acoustic Representation and Feature Selection}

Recordings were decoded as mono waveforms at 16~kHz, silence-trimmed at a 35-dB threshold, peak-normalized, and padded or cropped to fixed durations (4~s for cough, 5~s for breathing and speech). Quality screening required minimum duration (0.5~s cough, 1.0~s breathing/speech), maximum silence ratio (0.70), and maximum clipping ratio (0.01).

Three descriptor banks were merged: openSMILE ComParE 2016 functionals (6,373 features), INTERSPEECH 2010 paralinguistic functionals (1,586 features), and a project acoustic bank (2,177 features), yielding 10,140 numeric candidates \cite{eyben2010opensmile,schuller2010is10,schuller2013compare}.

\textbf{Critical:} Feature ranking was computed \emph{only} on training data (N=1,460). LightGBM gain importance was calculated separately for each modality, normalized within modality, and averaged across modalities. The 800 highest-scoring feature names were frozen; the same names were used in validation and test without recomputing ranks or scores.

\subsection{Classifiers and Modality Models}

Four classifier families were evaluated: LightGBM, XGBoost, CatBoost, and a sigmoid-calibrated radial-basis SVC \cite{ke2017lightgbm,chen2016xgboost,prokhorenkova2018catboost,cortes1995svm}. Hyperparameters were fixed before evaluation. SMOTE (3 nearest neighbors) was applied \emph{only} to training rows after feature filtering; validation and test distributions were never resampled \cite{chawla2002smote}.

Recording-level probabilities were averaged within participant and modality. For each modality, validation AUROC (with AUPRC tie-breaker) selected either one classifier or the mean of the four validation-ranked classifiers. This selected a four-model mean for cough and breathing, and LightGBM for speech.

\subsection{Fusion Rules and Threshold Selection}

Three fusion rules were evaluated using \emph{validation} participants only:

\textbf{Uniform mean:} Average of available modality probabilities.

\textbf{Validation-weighted:} Weights derived from validation AUPRC as $w_m = \max(\mathrm{AUPRC}_m - 0.5, 0.01) / \sum_j r_j$.

\textbf{Logistic stack:} Class-balanced L2 logistic regression fitted on aligned validation probabilities.

\textbf{Critical:} The primary fusion rule was uniform averaging. Validation AUROC selected only \emph{which modalities} to include in the mean. Test metrics were excluded from this selection. The balanced-accuracy threshold was selected on validation data.

\subsection{Baselines and Negative Controls}

\textbf{Metadata-only model:} A LightGBM model trained on available demographic and recording-protocol variables (age, gender, recording date, location, device) to establish non-audio predictiveness.

\textbf{Shuffled-label control:} The complete pipeline was rerun after randomly permuting training labels to establish the absence of gross leakage.

\subsection{Uncertainty Quantification}

Primary endpoint was participant-level AUROC. A nonparametric bootstrap resampled test participants 1,000 times to obtain 95\% confidence intervals. Paired DeLong tests compared AUROCs where the same participants had predictions from both models \cite{delong1988auc}.

\section{Results}

\subsection{Validation-Only Model Selection}

Among uniform-mean modality combinations evaluated on validation data, cough plus speech achieved AUROC 0.842 (AUPRC 0.800), ranking first among predefined combinations. The three-modality mean achieved validation AUROC 0.841. This validation-only rule identified cough plus speech as the primary configuration.

\subsection{Participant-Disjoint Test Performance}

Table~\ref{tab:results} reports test performance. Speech was the strongest single modality with AUROC 0.888 (95\% CI 0.842--0.930). The validation-selected cough--speech mean achieved AUROC 0.895 (95\% CI 0.852--0.933), AUPRC 0.862, balanced accuracy 0.808, and F1 0.729 on 314 test participants with both modalities available.

\begin{table}[!t]
\caption{Participant-Level Test Results with 95\% Bootstrap Confidence Intervals}
\label{tab:results}
\centering
\footnotesize
\setlength{\tabcolsep}{2.8pt}
\begin{tabular}{lrrrr}
\toprule
System & $n$ & AUROC (95\% CI) & AUPRC & F1 \\
\midrule
Breathing, top-4 mean & 312 & 0.828 (0.775--0.877) & 0.734 & 0.654 \\
Cough, top-4 mean     & 316 & 0.862 (0.812--0.908) & 0.804 & 0.705 \\
Speech, LightGBM      & 314 & 0.888 (0.842--0.930) & 0.833 & 0.740 \\
\midrule
\textbf{Cough+speech mean} & \textbf{314} & \textbf{0.895 (0.852--0.933)} & \textbf{0.862} & \textbf{0.729} \\
\bottomrule
\end{tabular}
\end{table}

The fusion estimate was 0.007 AUROC above speech alone. Paired DeLong test on the 314 participants with both predictions: difference = 0.007, 95\% CI $-0.021$ to 0.035, $p=0.62$. This is not statistically significant; the data do not establish fusion superiority.

\subsection{Baseline and Negative Controls}

The metadata-only model achieved test AUROC 0.964 (95\% CI 0.938--0.984). After full pipeline retraining with shuffled labels, mean AUROC was 0.503.

\subsection{Literature Context}

Table~\ref{tab:context} places the estimate in context of prior Coswara studies.

\begin{table}[!t]
\caption{Protocol-Aware Comparison with Prior Coswara Studies}
\label{tab:context}
\centering
\footnotesize
\setlength{\tabcolsep}{2.5pt}
\begin{tabular}{llll}
\toprule
Study & Inputs & Protocol & AUROC \\
\midrule
Bhattacharya \cite{bhattacharya2023coswara} & 9 sounds + symptoms & Internal 70/15/15 & 0.915 \\
Chetupalli \cite{chetupalli2023multimodal} & Audio fusion & Subject-disjoint & 0.880 \\
Truong \cite{truong2024fair} & 7 sounds & Fixed test & $0.866 \pm 0.012$ \\
\midrule
\textbf{This work} & Cough + speech & Participant-disjoint & \textbf{0.895} \\
\bottomrule
\end{tabular}
\par\vspace{2pt}\scriptsize 
\textbf{Note:} These are not head-to-head comparisons. Cohort snapshots, label definitions, modalities, and partition rules differ. They establish that the present estimate is within the range of prior internal Coswara results.
\end{table}

\section{Discussion}

\subsection{Principal Findings}

A multimodal audio model achieved AUROC 0.895 under rigorously isolated participant-disjoint evaluation, with 95\% CI 0.852--0.933. This is comparable to prior Coswara estimates (0.88--0.92) but does not establish superiority. The fusion gain over speech alone (0.007 AUROC) was not statistically significant.

Metadata alone achieved AUROC 0.964, substantially higher than audio. Shuffled-label retraining produced AUROC 0.503, establishing the absence of gross label leakage.

\subsection{Relation to Prior Evidence}

These results extend Han and Coppock's methodological findings. Han showed biased splits can inflate AUROC by 0.19 \cite{han2022realistic}. Coppock showed matching reduces audio AUROC from 0.85 to 0.62 and symptoms outperform audio \cite{coppock2024audio}. Our results show that under participant-disjoint, validation-only evaluation, audio can achieve strong internal discrimination, but metadata is more predictive.

This confirms that evaluation protocol is critical: internal performance does not establish external validity, temporal stability, or superiority over non-audio baselines.

\subsection{Limitations}

\textbf{First,} this is a retrospective analysis of crowdsourced data. Labels were self-reported, not PCR-confirmed by the present authors. The estimate applies to the Coswara label definition, not to clinical diagnosis.

\textbf{Second,} validation data served multiple selection roles: modality model, fusion combination, and threshold. The 800-feature budget was one of three explored during development. A fully nested outer evaluation would estimate the complete procedure more conservatively.

\textbf{Third,} only 314 test participants had both cough and speech predictions. The fusion estimate applies to this subset, not to all participants.

\textbf{Fourth,} this study evaluates internal performance only. Temporal drift, cross-dataset transfer, and prospective validation require separate protocols and are outside the present scope.

\textbf{Fifth,} COUGHVID external transfer was not tested in this conference paper. The full reliability analysis (temporal, external, mechanism) is reserved for a journal extension.

\subsection{Implications}

Strong internal discrimination is achievable under rigorous evaluation. However, metadata outperforming audio raises questions about what the audio model learns. Future work should test whether audio adds information beyond metadata on matched participants, and whether any internal estimate predicts external transfer.

\section{Conclusion}

Under participant-disjoint evaluation with validation-only model selection, a Coswara multimodal model achieved AUROC 0.895 (95\% CI 0.852--0.933). This is comparable to prior internal estimates. The fusion gain over speech alone was not statistically significant. Metadata alone achieved higher discrimination (AUROC 0.964). The estimate applies only to the Coswara test cohort and does not establish external validity, temporal stability, or clinical readiness.

\section*{Data Availability}

This work uses public Coswara data under its access terms \cite{bhattacharya2023coswara}. Analysis code and frozen configurations will accompany the archived release.

\section*{Ethics Statement}

This is secondary analysis of public data. No new participants were recruited. Consent and collection procedures are described in the source publication.

\bibliographystyle{IEEEtran}
\bibliography{references}

\end{document}